\begin{lemma}{Dvoretzky--Rogers--Lemma}{DvoretzkyRogersLemma}
Seien $n, m \in \N$ mit $n \ge 4m^2$, sowie $(E, \norm\cdot_E)$ ein normierter
$\R$--Vektorraum mit $\dim E = n$.\\
Dann gibt es $x \in E^m$ mit
\[\frac 1 2 \norm\cdot_{\infty, m} \le \norm\cdot_E \circ \bprod \cdot x_m
\le \norm\cdot_{2,m}\]
\end{lemma}
\begin{proof}
Ein Beweis dieser Aussage befindet sich in \cite{Kad}, dort Lemma 6.2.1.
Leider enth�lt die dortige Absch�tzung in der letzten Zeile einen kleinen
Fehler, der aber wie folgt (gem�� dem urspr�nglichen Beweis von Dvoretzky und
Rogers, vgl. \cite{DR}) behoben werden kann (die Notation lehnt sich an jene
aus \cite{Kad} an):
\begin{quote}
Sei $i \in \haken m$ mit $i+1 \in \haken m$ und die Behauptung gelte f�r alle
$k \le i$.\\
Wegen $i < i+1 \le m \le \frac {\sqrt n} 2$ gilt:
\[i\sqrt n + i^2 + i = i\left(\sqrt n+i+1\right)
\le\frac{\sqrt n}2\left(\sqrt n+\frac{\sqrt n}2\right)
= \frac 3 4 n\text.\]
Daher ist
\[4\frac i {\sqrt n} + 4\frac{i^2}n+ 4 \frac i n \le 3,\]
also
\begin{align}\left(1+\frac {2i}{\sqrt n} \right)^2
&= 1 + 4 \frac i{\sqrt n} +  4\frac {i^2}n \le 4 - 4\frac i n 
= 4 \left(1- \frac i n \right)\text. \tag{$\ast$}\label{drleq:1}
\end{align}
Damit gilt wie im Beweis angegeben:
\begin{align*}
\abs{t_{i+1}}^2 &\le \frac n {n-i}\left(1+ 2\sqrt{i}\sqrt{\frac i n} \right)^2
=\frac 1{1-\frac in}\left(1+\frac{2i}{\sqrt n}\right)^2
\overset{ \eqref{drleq:1}}\le 4\text.
\end{align*}
\end{quote}
Dies liefert genau die f�r den Beweis ben�tigte Aussage.
\end{proof}